\documentclass[tikz,border=5pt,12pt]{standalone}
\usepackage{resistencia}
\begin{document}
\begin{tikzpicture}[x=1cm,y=1cm,etq/.style={etiq,fill=white,inner sep=1.5pt}]
% ---- eje (vista lateral) ----
\def\R{0.9}\def\L{6.2}
\filldraw[cuerpo] (0,\R) -- (\L,\R) arc[start angle=90,end angle=-90,x radius=0.32,y radius=\R] -- (0,-\R) arc[start angle=-90,end angle=-270,x radius=0.32,y radius=\R] -- cycle;
\draw[contorno] (\L,0) ellipse (0.32 and \R);
\draw[eje,dash pattern=on 6pt off 3pt,thin] (-0.9,0) -- (\L+1.2,0);
% pares en los extremos
\draw[fuerza,<-] (-0.95,0.9) arc[start angle=100,end angle=-100,x radius=0.35,y radius=0.95];
\node[etq,text=cobre] at (-1.35,0) {$T$};
\draw[fuerza] (\L+0.95,0.9) arc[start angle=80,end angle=-80,x radius=0.35,y radius=0.95];
\node[etq,text=cobre] at (\L+1.4,0) {$T$};
% plano de corte
\draw[corte] (3.3,-1.5) -- (3.3,1.5);
\node[etq,text=violeta] at (3.3,1.85) {plano de corte $\Pi$};
% ---- sección transversal ampliada ----
\begin{scope}[xshift=10.6cm]
  \filldraw[cuerpo] (0,0) circle (1.7);
  \fill (0,0) circle (1.6pt);
  \foreach \r in {0.55,1.1,1.65}{
    \draw[gris] (0,0) circle (\r);
    \foreach \a in {0,90,180,270}{
      \pgfmathsetmacro{\px}{\r*cos(\a)}\pgfmathsetmacro{\py}{\r*sin(\a)}
      \pgfmathsetmacro{\dx}{-0.38*\r*sin(\a)}\pgfmathsetmacro{\dy}{0.38*\r*cos(\a)}
      \draw[fuerza] (\px,\py) -- ++(\dx,\dy);
    }
  }
  \node[etq,text=cobre,anchor=west] at (2.0,0.3) {$\tau_{\max}$};
  \node[etq,anchor=west] at (2.0,-0.3) {en la superficie};
  \node[etq,text width=5.6cm,align=center] at (0.3,-2.75) {$\tau=\dfrac{T\,r}{J}$: crece con el radio $r$};
  \node[etq,text=violeta] at (0,2.25) {sección transversal};
\end{scope}
\end{tikzpicture}
\end{document}
